\section{Phase selection on a fixed Hermitian--Einstein bundle}
\label{sec:parallel-phase}

The component metrics required for spatial initial data are selected by
an intrinsic fixed-background argument. Let $(Y,\omega)$ be a connected
compact K\"ahler manifold of volume $V$. Let $(E,h_0)$ be a
Hermitian--Einstein bundle of rank $r$, with Chern connection $D$ and
$i\Lambda F_D=\lambda I$. Write $h=h_0H$ and use the normalized equation
\begin{equation}\label{eq:fixed-flow}
 H^{-1}H_s=-2(i\Lambda F_h-\lambda I),\qquad \det H=1.
\end{equation}
The base and connection are fixed throughout this section. The positive
gap used here is not asserted to be uniform on a degenerating family.

Let $\mathcal A$ be the algebra of $D$-parallel complex endomorphisms,
and let $\Pi$ be the $L^2$ orthogonal projection onto $\mathcal A$.
Denote by $\lambda_\perp>0$ the first positive eigenvalue of $D^*D$
on self-adjoint endomorphisms, with its parallel kernel omitted.

\begin{theorem}\label{thm:parallel-phase}
Suppose $H(0)$ is smooth, $\det H(0)=1$ and $cI\le H(0)\le CI$,
where $0<c\le C<\infty$. There is a unique positive parallel
endomorphism $P=P(H(0))$, with $\det P=1$, to which $H(s)$ converges
in $L^2$ and uniformly. Put
\begin{equation}\label{eq:phase-functionals}
 M=\Pi H,\qquad \mathcal J=\log\det M,\qquad
 W=V^{-1}\int_Y|H-M|^2\,dV,\qquad
 \gamma=\frac{2c^2\lambda_\perp}{C^2}.
\end{equation}
Then $M$ decreases in matrix order to $P$, and
\begin{equation}\label{eq:phase-estimates}
 \frac{W}{2C^2}\le \mathcal J\le\frac{W}{2c^2},\qquad
 \mathcal J(s)\le \mathcal J(0)e^{-\gamma s},\qquad
 \tr(M(s)-P)\le C\mathcal J(s).
\end{equation}
In particular,
\begin{equation}\label{eq:phase-l2}
 V^{-1}\int_Y|H(s)-P|^2\,dV
       \le C^2(2\mathcal J(0)+\mathcal J(0)^2)e^{-\gamma s}.
\end{equation}
For $s\ge1$ the squared affine distance satisfies
\begin{equation}\label{eq:phase-uniform}
 \sup_Y d_r(H(s),P)^2
       \le K_Y(2\mathcal J(0)+\mathcal J(0)^2)e^{-\gamma(s-1)}.
\end{equation}
Here $K_Y$ also depends on $c,C,r$ and the fixed background. For two
initial metrics in a common positive range, the selected phases obey
\begin{equation}\label{eq:phase-contraction}
 V\sigma_r(P(H(0)),P(K(0)))
       \le\int_Y\sigma_r(H(0),K(0))\,dV.
\end{equation}
Consequently phase selection is Lipschitz on such a range. If
$\mathcal A=\CC I$, the selected phase is $I$.
\end{theorem}

\begin{proof}
Decompose the unitary holonomy representation at a base point into
inequivalent irreducible representations and their multiplicity spaces.
Parallel transport of the invariant summands gives a global orthogonal
parallel decomposition
\begin{equation}\label{eq:holonomy-splitting}
 E=\bigoplus_a F_a\otimes\CC^{m_a},\qquad
 D=\bigoplus_a D_a\otimes I,
 \qquad \mathcal A=\bigoplus_a I_{F_a}\otimes\End(\CC^{m_a}).
\end{equation}
The transported subspaces are independent of the chosen path because
they are holonomy invariant. This does not require individual parallel
vectors on $F_a$. The splitting is also holomorphic since the
$(0,1)$ part of $D$ is the holomorphic structure.

If $r_a=\operatorname{rank}F_a$ and $B_{aa}$ is a diagonal block of
an endomorphism section, direct testing against every parallel
endomorphism gives
\begin{equation}\label{eq:parallel-projection}
 \Pi B=\bigoplus_a I_{F_a}\otimes
       \frac1{r_aV}\int_Y\tr_{F_a}B_{aa}\,dV.
\end{equation}
The partial trace is globally defined using the fixed metric on $F_a$.
Diagonal restriction, partial trace and integration preserve positivity.
Thus $\Pi$ is positive, unital and adjoint preserving. It preserves the
normalized integrated trace, and, for any parallel endomorphism $B$,
\begin{equation}\label{eq:projection-trace}
 \tr(B\Pi Q)=V^{-1}\int_Y\tr(BQ)\,dV.
\end{equation}
Order comparison for \eqref{eq:fixed-flow} preserves $cI\le H\le CI$;
its scalar determinant equation preserves $\det H=1$. Hence
$cI\le M\le CI$.

We need the full relative metric equation. In unitary normal base
coordinates with $\omega=(i/2)\sum dz_\alpha\wedge d\bar z_\alpha$
at the point, expansion of $\bar\partial_D(H^{-1}\partial_DH)$ gives
\begin{equation}\label{eq:relative-fixed-equation}
 H_s=-D^*DH-Q(H),\qquad
 Q(H)=4\sum_\alpha(D_\alpha H)^*H^{-1}(D_\alpha H)\succeq0.
\end{equation}
The second-order term has scalar principal part
$4\sum\partial_\alpha\partial_{\bar\alpha}H$. Its difference from
$-D^*DH$ on endomorphism-valued zero-forms is the contracted curvature
commutator, which vanishes because $i\Lambda F_D=\lambda I$.
Differentiating $H^{-1}$ gives precisely the displayed quadratic term.
In particular, it cannot be dropped for matrix-valued $H$.
Since $\Pi D^*DH=0$, positivity of the projection proves
\begin{equation}\label{eq:monotone-mean}
 M'=-\Pi Q(H)\preceq0.
\end{equation}
It follows that $M$ has a parallel positive limit $P$ between $cI$ and $CI$.

Taylor's integral remainder for $\log\det$ on the segment from $M$
to $H$ has Hessian between $-c^{-2}$ and $-C^{-2}$ in Frobenius norm.
The integrated linear term vanishes, since
\[
 V^{-1}\int_Y\tr(M^{-1}(H-M))\,dV=0
\]
by \eqref{eq:projection-trace}. Integrating $\log\det H=0$ gives
the two bounds for $\mathcal J$ in \eqref{eq:phase-estimates}.
Next, the same trace identity and positivity give
\begin{align*}
 -\mathcal J'&=V^{-1}\int_Y\tr(M^{-1}Q(H))\,dV\\
     &\ge C^{-2}V^{-1}\int_Y|DH|^2\,dV
      \ge C^{-2}\lambda_\perp W\ge\gamma \mathcal J.
\end{align*}
For the first inequality, remove both positive matrix weights using
their lower bounds. Only then use
$4\sum_\alpha\tr((D_\alpha H)^*D_\alpha H)=|DH|^2$.
The gap applies because $H-M$ is orthogonal to the whole parallel
algebra and $D(H-M)=DH$. This proves exponential decay of $\mathcal J$.

By continuity, $\log\det P=0$. Integrating the derivative of
$\log\det$ on the positive segment from $P$ to $M$ gives
$\tr(M-P)\le C\mathcal J$. Orthogonality to parallel sections gives the exact
decomposition
\[
 V^{-1}\int_Y|H-P|^2\,dV=W+|M-P|^2
       \le 2C^2\mathcal J+C^2\mathcal J^2.
\]
Here $M-P\succeq0$, so its norm is at most its trace. This proves
\eqref{eq:phase-l2}.

The metric $h_0P$ has the same Chern connection as $h_0$ and is
stationary. The unforced trace comparison from
\eqref{eq:forced-trace}, followed by the time-one scalar heat kernel
on the fixed compact base, bounds the supremum of
$\sigma_r(H(s),P)$ by a fixed multiple of its integral at $s-1$.
Equivalence of trace distance and squared norm on the fixed range
proves \eqref{eq:phase-uniform}. This uses no derivative bound on
arbitrary bounded initial metrics and no degenerating heat kernel.

Finally, integrate the trace comparison for two exact flows. Uniform
convergence to their parallel limits proves \eqref{eq:phase-contraction}.
It also gives uniqueness, and the equivalence of the distances gives
the stated Lipschitz dependence. A positive scalar endomorphism with
determinant one is the identity.
\end{proof}

On a bundle with repeated irreducible factors the phase belongs to the
actual algebra \eqref{eq:holonomy-splitting}; it is not a constant matrix
in a chosen local frame on the whole bundle. In the next section the
fixed component bundles supply this algebra and its positive gap.
The global degeneration estimate will be established separately.
